\documentclass[10pt,a4paper]{amsart}
\usepackage[margin=19mm]{geometry}
\usepackage{amsmath,amssymb,mathtools}
\usepackage[scr=rsfs,cal=euler]{mathalfa}
\usepackage{xfrac}
\usepackage[unicode,hidelinks,bookmarks=false]{hyperref}
\newcommand{\ie}{i.e.\ }
\newcommand{\cf}{cf.\ }
\newcommand{\resp}{respectively\ }
\newcommand{\Subsection}{\subsection}
\providecommand{\nobreakdash}{\nobreak\hbox{-}}
\setlength{\emergencystretch}{2em}
\multlinegap0pt
\usepackage{bm}
\usepackage{bbm}
\usepackage{dsfont}
\usepackage{xspace}
\usepackage{cancel}
\usepackage{mathtools}
\usepackage{amsthm}
\makeatletter
\@ifundefined{lemma}{\newtheorem{lemma}{Lemma}}{}
\makeatother
\makeatletter
\expandafter\ifx\csname cleverproof\endcsname\relax
\newenvironment{cleverproof}[1]
  {\begin{proof}[Proof of \Cref{#1}]}
  {\end{proof}}
\fi
\makeatother
\title[Bayesian Optimal Stopping with Maximum Value Knowledge]{Bayesian Optimal Stopping with Maximum Value Knowledge}
\author{Pieter Kleer, Daan Noordenbos}
\date{Source: arXiv:2507.03497v1 (CC BY 4.0)}
\begin{document}
\maketitle

\noindent\emph{Source and licence.} Bayesian Optimal Stopping with Maximum Value Knowledge. Version arXiv:2507.03497v1 (CC BY 4.0), distributed under the Creative Commons Attribution 4.0 International licence, as recorded in the arXiv licence metadata for this exact version and shown on the abstract page https://arxiv.org/abs/2507.03497v1. Full rights notice: licenses/arxiv-2507.03497-CC-BY-4.0.txt.\par\smallskip

\noindent\textit{Editorial scope.} Counted unit: the Lemma whose source label is \texttt{lemma:technical-upperbound-argument} in the article (source0.tex lines 1385--1390, source label \texttt{lemma:technical-upperbound-argument}), delivered together with the article's own complete original proof of it (source0.tex lines 1391--1471, one \texttt{proof} environment of 81 lines). No mathematics is written, repaired or re-derived: statement and proof are the article's own text line for line. Required context retained uncounted: lemma:delta-choice (lines 858--862); lemma:stopping-time-difference-equation-1 (lines 1274--1280); lemma:stopping-time-difference-equation-2 (lines 1365--1371). Every definition, display and label of those lines is retained unchanged. Omitted from this package and not redistributed: 1--857, 863--1273, 1281--1364, 1372--1384, 1472--1476. Nothing inside a retained excerpt is omitted or altered: every retained body is delivered exactly as the article's lines are written, so no omission receipt of any kind accompanies this package. Citations: the retained excerpts carry no citation command at all, so no bibliography entry of the article is transcribed and no reference is redistributed. Reference adaptation: none was needed and none was made; every reference inside the delivered unit resolves inside it, so no unresolved reference or citation is produced. Printed slips kept literal: no printed word, symbol, hypothesis or number of the counted statement or of its proof was altered. Layout and class: the article's own class is \texttt{article} with packages including geometry, soul, preamble, cleveref, booktabs; the wrapper is the lane's pinned \texttt{amsart} editorial wrapper at 10pt on A4, which restates the theorem-like environments the retained text needs and adds the article's own macro definitions that the retained text needs (none), with extra wrapper packages (bm, bbm, dsfont, xspace, cancel, mathtools). The delivered numbering is the unit's own. Assets: no retained span contains a figure environment, a graphics-inclusion command, a PDF-page inclusion, a table or a TikZ picture; no image file is copied into this package, the package declares no asset dependency and no image file is redistributed with it. Licence: version arXiv:2507.03497v1 of this article is distributed under the Creative Commons Attribution 4.0 International licence, as recorded in the arXiv licence metadata for this exact version and shown on the abstract page https://arxiv.org/abs/2507.03497v1. A full rights notice is delivered at licenses/arxiv-2507.03497-CC-BY-4.0.txt. Characters: every retained excerpt is pure ASCII, so the delivered engine, XeLaTeX with its default Latin Modern text font, reports no missing character.
    \begin{lemma}
        \label{lemma:delta-choice}
        Let $M(t;\epsilon):=\frac{\Pi^*+\epsilon}{\mathbb{P}(X_{\max}\ge t)}$, $f(t;\epsilon)=M(t;\epsilon)-t$ and $\bar{p}(\epsilon)$ the minimiser of $f$. 
        If $F_{\max}$ is continuously differentiable at $p^*$, then there exist $\epsilon_0'>0$ and $\delta>0$ (independent of $\epsilon$) such that for $\epsilon<\epsilon_0'$ the following holds: $M$ is increasing and continuous on $(p^*-\delta,p^*+\delta)$, and $f$ is non-negative, decreasing on $(p^*-\delta,\bar{p}(\epsilon)]$ and increasing on $[\bar{p}(\epsilon),p^*+\delta)$
    \end{lemma}

\begin{lemma}
    \label{lemma:stopping-time-difference-equation-1}
    Let $I$ be a continuous and increasing function which can be written as $I(t)=t+f(t)$ where $f$ is a decreasing and non-negative function on $[a,b]$ and let $I^k=I\circ I^{k-1}$, then
    $$
    \inf\{k\in\mathbb{N} : I^k(a)\ge b\}\le 1 + \int_a^bf(t)^{-1}\mathrm{d}t.
    $$
\end{lemma}

\begin{lemma}
    \label{lemma:stopping-time-difference-equation-2}
    Let $I$ be a continuous and increasing function which can be written as $I(t)=t+f(t)$ where $f$ is an increasing and non-negative function on $[a,b]$. Let $I^k=I\circ I^{k-1}$ and $\tilde{f}(t)=t-I^{-1}(t)$, then
    $$
    \inf\{k\in\mathbb{N}  : I^k(a)\ge b\}\le 1 + \int_a^b\tilde{f}(t)^{-1}\mathrm{d}t.
    $$
\end{lemma}

\begin{lemma}\label{lemma:technical-upperbound-argument}
Let $\Pi^*=\Pi^*(F_{\max})$, $M(t;\epsilon):=\frac{\Pi^*+\epsilon}{\mathbb{P}(X_{\max}\ge t)}$, $M^k=M\circ M^{k-1}$, $k(a,b;\epsilon)=\inf\{k:M^k(a;\epsilon)\ge b\}$ and $\delta$ chosen according to Lemma \ref{lemma:delta-choice}, then 
    $$
    \lim_{\epsilon\to 0}\epsilon^{\frac{1}{2}}k(p^*-\delta,p^*+\delta;\epsilon)\le\sqrt{2\pi^2\frac{(1-F_{\max}(p^*))^2}{p^*F''_{\max}(p^*)+2F'_{\max}(p^*)}}.
    $$
\end{lemma}

\begin{proof}
    Let $f(t;\epsilon)=M(t;\epsilon)-t$ and $\bar{p}$ be the minimiser of $f$. By the triangle like inequality,
    $$
    k(p^*-\delta,p^*+\delta;\epsilon)\le k(p^*-\delta,\bar{p};\epsilon)+k(\bar{p},p^*+\delta;\epsilon).
    $$
    Choosing $\delta$ according to Lemma \ref{lemma:delta-choice} allows us to apply Lemmas \ref{lemma:stopping-time-difference-equation-1} and \ref{lemma:stopping-time-difference-equation-2}, which give us the following upper bounds on the above two terms
    \begin{align}
        k(p^*-\delta,\bar{p};\epsilon)&\le 1+\int_{p^*-\delta}^{\bar{p}}f(t;\epsilon)^{-1}\mathrm{d}t,\label{eq:bound-left}\\
        k(\bar{p},p^*+\delta;\epsilon)&\le1+\int_{\bar{p}}^{p^*+\delta}\left(t-M^{-1}(t;\epsilon)\right)^{-1}\mathrm{d}t.\label{eq:bound-right}
    \end{align}
    $f(t;\epsilon)$ is known and $t-M^{-1}(t;\epsilon)=t-F^{-1}_{\max}\left(1-\frac{\Pi^*+\epsilon}{t}\right)$.
    We now proceed with the limit argument.    
    Let $\zeta_1(t;\epsilon)=f(t;\epsilon)$ and $\zeta_2(t;\epsilon)=t-M^{-1}(t;\epsilon)$. For $|s|\le \delta$ we have by Taylor's theorem that,
    $$
    \zeta_i(p^*+s)\ge\xi_i(p^*+s;\epsilon):=\zeta_i(p^*)+\zeta_i'(p^*)s+\frac{1}{2}\zeta_i''(p^*)s^2-L|s|^3, 
    $$
    where $L=\max\{L_1,L_2\}$ and $L_i=\sup\{|\frac{1}{6}\zeta_i'''(p^*+s)|:|s|\le\delta\}$. The $\xi_i$'s can be written as
    $$
    \xi_i(p^*+s;\epsilon):=a_i+b_is+c_is^2-L|s|^3,
    $$
    where
    \begin{align*}    
        a_1&=\frac{\Pi^*+\epsilon}{1-F_{\max}(p^*)}-p^*=\frac{\epsilon}{1-F_{\max}(p^*)},\\
        a_2&=p^*-F^{-1}_{\max}\left(1-\frac{\Pi^*+\epsilon}{p^*}\right),\\
        b_1&=\frac{(\Pi^*+\epsilon)F'_{\max}(p^*)}{(1-F_{\max}(p^*))^2}-1,\\
        b_2&=1-\frac{\Pi^*+\epsilon}{(p^*)^2F'_{\max}\left(F^{-1}_{\max}\left(1-\frac{\Pi^*+\epsilon}{p^*}\right)\right)},\\
        c_1&=(\Pi^*+\epsilon)\frac{F''_{\max}(p^*)(1-F_{\max}(p^*))+2F'_{\max}(p^*)^2}{2(1-F_{\max}(p^*))^3},\\
        c_2&=\frac{(\Pi^*+\epsilon)}{(p^*)^3F'_{\max}\left(F^{-1}_{\max}\left(1-\frac{\Pi^*+\epsilon}{p^*}\right)\right)}+\frac{(\Pi^*+\epsilon)^2F''_{\max}\left(F^{-1}_{\max}\left(1-\frac{\Pi^*+\epsilon}{p^*}\right)\right)}{2(p^*)^4F'_{\max}\left(F^{-1}_{\max}\left(1-\frac{\Pi^*+\epsilon}{p^*}\right)\right)^3}.
    \end{align*}
    We will now show that $a_i=a\epsilon+O(\epsilon^2)$, $b_i=O(\epsilon)$ and $c_i=c+O(\epsilon)$ where $a=(1-F_{\max}(p^*))^{-1}$ and $c=\frac{p^*F''_{\max}(p^*)+2F'_{\max}(p^*)}{2(1-F_{\max}(p^*))}$. 
    
    Recall that $1-F_{\max}(p^*)=p^*F'_{\max}(p^*)$, since $p^*$ solves the first order condition.
    Firstly, $a_1=a\epsilon$ and $a_2=a\epsilon+O(\epsilon^2)$ follows from the fact that $a_2=0$ for $\epsilon=0$ and that 
    \begin{align*}
        \left. \frac{\mathrm{d} a_2}{\mathrm{d} \epsilon} \right|_{\epsilon = 0}&=
    \left. \left(p^*F'_{\max}\left(F^{-1}_{\max}\left(1-\frac{\Pi^*+\epsilon}{p^*}\right)\right)\right)^{-1} \right|_{\epsilon = 0}\\
    &=(p^*F'_{\max}(p^*))^{-1}=(1-F_{\max}(p^*))^{-1}.
    \end{align*}
    Secondly, for $\epsilon=0$ 
    \begin{align*}    
        b_1&=\frac{(\Pi^*)F'_{\max}(p^*)}{(1-F_{\max}(p^*))^2}-1=\frac{p^*F'_{\max}\left(p^*\right)}{1-F_{\max}(p^*)}-1=0,\\
        b_2&=1-\frac{\Pi^*}{(p^*)^2F'_{\max}\left(F^{-1}_{\max}\left(1-\frac{\Pi^*}{p^*}\right)\right)}=1-\frac{1-F_{\max}(p^*)}{p^*F'_{\max}\left(p^*\right)}=0.
    \end{align*}
    Therefore $b_1$ and $b_2$ are $O(\epsilon)$. Lastly, for $\epsilon=0$ 
    \begin{align*}    
        c_1&=(\Pi^*)\frac{F''_{\max}(p^*)(1-F_{\max}(p^*))+2F'_{\max}(p^*)^2}{2(1-F_{\max}(p^*))^3}\\
        &=\frac{p^*F''_{\max}(p^*)(1-F_{\max}(p^*))+2p^*F'_{\max}(p^*)^2}{2(1-F_{\max}(p^*))^2}\\
        &=\frac{p^*F''_{\max}(p^*)+2F'_{\max}(p^*)}{2(1-F_{\max}(p^*))},\\
        c_2&=\frac{\Pi^*}{(p^*)^3F'_{\max}\left(F^{-1}_{\max}\left(1-\frac{\Pi^*}{p^*}\right)\right)}+\frac{(\Pi^*)^2F''_{\max}\left(F^{-1}_{\max}\left(1-\frac{\Pi^*}{p^*}\right)\right)}{2(p^*)^4F'_{\max}\left(F^{-1}_{\max}\left(1-\frac{\Pi^*}{p^*}\right)\right)^3}\\
        &=\frac{1-F_{\max}(p^*)}{(p^*)^2F'_{\max}\left(p^*\right)}+\frac{(1-F_{\max}(p^*))^2F''_{\max}\left(p^*\right)}{2(p^*)^2F'_{\max}\left(p^*\right)^3}\\
        &=\frac{1}{p^*}+\frac{F''_{\max}\left(p^*\right)}{2F'_{\max}\left(p^*\right)}=\frac{p^*F''_{\max}(p^*)+2F'_{\max}(p^*)}{2p^*F'_{\max}(p^*)}=\frac{p^*F''_{\max}(p^*)+2F'_{\max}(p^*)}{2(1-F_{\max}(p^*))}.
    \end{align*}
    Therefore, $c_i=c+O(\epsilon)$. From (\ref{eq:bound-left}) and (\ref{eq:bound-right}) and the fact that $\zeta_i$ is lower bounded by $\xi_i$ we obtain that
    \begin{align*}
        \lim_{\epsilon\to 0}\epsilon^{\frac{1}{2}}\left(k(p^*-\delta,\bar{p};\epsilon)+k(\bar{p},p^*+\delta;\epsilon)\right)&\le\lim_{\epsilon\to 0}\epsilon^{\frac{1}{2}}\left(\int_{p^*-\delta}^{\bar{p}}\frac{\mathrm{d}t}{\zeta_1(t;\epsilon)}+\int_{\bar{p}}^{p^*+\delta}\frac{\mathrm{d}t}{\zeta_2(t;\epsilon)}\right)\\
        &\le\lim_{\epsilon\to 0}\epsilon^{\frac{1}{2}}\left(\int_{p^*-\delta}^{\bar{p}}\frac{\mathrm{d}t}{\xi_1(t;\epsilon)}+\int^{p^*+\delta}_{\bar{p}}\frac{\mathrm{d}t}{\xi_2(t;\epsilon)}\right).
    \end{align*}
    The value of the right-hand side is driven by the first order dependence of $a_i,b_i$ and $c_i$ on $\epsilon$. Since these are the same for $i=1$ and $i=2$, we have that for $\xi_0(p^*+s;\epsilon):=a_0+b_0s+c_0s^2-L|s|^3$ with $a_0=a\epsilon+O(\epsilon^2)$, $b_0=O(\epsilon)$ and $c_0=c+O(\epsilon)$ that 
    \begin{align}
        \lim_{\epsilon\to 0}\epsilon^{\frac{1}{2}}\left(\int_{p^*-\delta}^{\bar{p}}\frac{\mathrm{d}t}{\xi_1(t;\epsilon)}+\int^{p^*+\delta}_{\bar{p}}\frac{\mathrm{d}t}{\xi_2(t;\epsilon)}\right)=\lim_{\epsilon\to 0}\epsilon^{\frac{1}{2}}\int_{p^*-\delta}^{p^*+\delta}\frac{\mathrm{d}t}{\xi_0(t;\epsilon)}.\label{eq:xi-0}
    \end{align}
    The function $\xi_0$ has roots above and below $p^*$, so $\delta$ needs to be small enough such that there are no roots in the interval $(p^*-\delta,p^*+\delta)$. To do so we make use of two facts. Firstly, as $\epsilon$ decreases so does $\xi_0$, resulting in roots closer to $p^*$. For $\epsilon=0$, $\xi_0$ has roots at $p^*\pm\frac{c}{L}$, so $\delta <\frac{c}{L}$ is necessary. Moreover, for $\epsilon=0$, $\xi_0$ has local maxima at $p^*\pm\frac{2}{3}\frac{c}{L}$, so we pick $\delta\le\frac{2}{3}\frac{c}{L}$. 
    Such a $\delta$ exists, since $L$ is continuous in $\delta$ at zero, because $F_{\max}$ is three times continuously differentiable at $p^*$. This is not in conflict with Lemma \ref{lemma:delta-choice}, since $\delta$ can always be chosen smaller.
    
    Denote the unique real root of $a_0+b_0t+c_0t^2+Lt^3=0$ by $r_1$ and the unique real root of $a_0+b_0t+c_0t^2-Lt^3=0$ by $r_2$. We have that $r_1=-\frac{c}{L}+O(\epsilon)$ and $r_2=\frac{c}{L}+O(\epsilon)$.   
    Partial fraction decomposition yields that
    $$
    \frac{1}{\xi_0(p^*+s;\epsilon)}=\begin{cases}
        \lambda_1\left(\frac{1}{s-r_1}-\frac{Ls+c_0+2r_1L}{Ls^{2}+\left(c_0+r_1L\right)s-a_0/r_1}\right),&\text{if } s\le 0,\\
        \lambda_2\left(\frac{1}{s-r_{2}}-\frac{Ls+2r_{2}L-c_0}{Ls^{2}+\left(r_{2}L-c_0\right)s+a_0/r_{2}}\right),&\text{if } s> 0,
    \end{cases}
    $$
    where $\lambda_1=\frac{r_1}{r_1^{2}\left(c_0+2r_1L\right)-a_0}$ and $\lambda_2=\frac{r_{2}}{r_{2}^{2}\left(c_0-2r_{2}L\right)-a_0}$. (\ref{eq:xi-0}) can now be evaluated,
    \begin{align*}
        \int_{p^*-\delta}^{p^*+\delta}\frac{\mathrm{d}s}{\xi_0(s;\epsilon)}=&\lambda_1\int_{-\delta}^0\left(\frac{1}{s-r_1}-\frac{Ls+c_0+2r_1L}{Ls^{2}+\left(c_0+r_1L\right)s-a_0/r_1}\right)\mathrm{d}s\\
        +&\lambda_2\int_0^{\delta}\left(\frac{1}{s-r_{2}}-\frac{Ls+2r_{2}L-c_0}{Ls^{2}+\left(r_{2}L-c_0\right)s+a_0/r_{2}}\right)\mathrm{d}s\\
        =&\lambda_1\log\left|r_{1}\right|-\lambda_1\log\left|r_{1}+\delta\right|-\lambda_1\frac{1}{2}\log\left|\frac{a_{0}}{r_{1}}\right|\\-&\lambda_1\frac{2(c_{0}+2r_{1}L)-\left(c_{0}+r_{1}L\right)}{\sqrt{-4L\frac{a_{0}}{r_{1}}-\left(c_{0}+r_{1}L\right)^{2}}}\arctan\left(\frac{\left(c_{0}+r_{1}L\right)}{\sqrt{-4L\frac{a_{0}}{r_{1}}-\left(c_{0}+r_{1}L\right)^{2}}}\right)\\+&\lambda_1\frac{1}{2}\log\left|L\delta^{2}-\left(c_{0}+r_{1}L\right)\delta-\frac{a_{0}}{r_{1}}\right|\\+&\lambda_1\frac{2(c_{0}+2r_{1}L)-\left(c_{0}+r_{1}L\right)}{\sqrt{-4L\frac{a_{0}}{r_{1}}-\left(c_{0}+r_{1}L\right)^{2}}}\arctan\left(\frac{-2L\delta+\left(c_{0}+r_{1}L\right)}{\sqrt{-4L\frac{a_{0}}{r_{1}}-\left(c_{0}+r_{1}L\right)^{2}}}\right)\\
        +&\lambda_2\left(\log\left|\delta-r_{2}\right|-\log\left|r_{2}\right|+\frac{1}{2}\log\left|\frac{a_{0}}{r_{2}}\right|\right)\\+&\lambda_2\frac{2\left(2r_{2}L-c_{0}\right)-\left(r_{2}L-c_{0}\right)}{\sqrt{4L\frac{a_{0}}{r_{2}}-\left(r_{2}L-c_{0}\right)^{2}}}\arctan\left(\frac{\left(r_{2}L-c_{0}\right)}{\sqrt{4L\frac{a_{0}}{r_{2}}-\left(r_{2}L-c_{0}\right)^{2}}}\right)\\-&\lambda_2\frac{1}{2}\log\left|L\delta^{2}+\left(r_{2}L-c_{0}\right)\delta+\frac{a_{0}}{r_{2}}\right|\\-&\lambda_2\frac{2\left(2r_{2}L-c_{0}\right)-\left(r_{2}L-c_{0}\right)}{\sqrt{4L\frac{a_{0}}{r_{2}}-\left(r_{2}L-c_{0}\right)^{2}}}\arctan\left(\frac{2L\delta+\left(r_{2}L-c_{0}\right)}{\sqrt{4L\frac{a_{0}}{r_{2}}-\left(r_{2}L-c_{0}\right)^{2}}}\right)
    \end{align*}
    Every term except the fourth and eighth vanishes when we multiply with $\epsilon^{\frac{1}{2}}$ and take the limit. The arctans converge to $\pm\pi$, as their arguments diverge, leaving only simple terms. Rudimentary calculus reveals that $\epsilon^{\frac{1}{2}}\int_{p^*-\delta}^{p^*+\delta}\frac{\mathrm{d}s}{\xi_0(s;\epsilon)}$ converges to $\pi(ac)^{-\frac{1}{2}}$. Substituting the expressions of $a$ and $c$ yields the result.
\end{proof}

\end{document}
